\section{7. Appendix A. Infographic
Prompts}\label{appendix-a.-infographic-prompts}

This appendix records one prompt per infographic so the visual layer can
be regenerated, refined, or re-styled without losing the design intent
or repo provenance.

\subsection{F01. Repository at a
glance}\label{f01.-repository-at-a-glance}

This prompt defines the root-map figure that lets a reader grasp the
repo's shape without reading the whole index first.

\begin{itemize}
\tightlist
\item
  Figure file: \texttt{f01\_repo\_at\_a\_glance.png}
\item
  Placement:
  \texttt{\#\#\#\ 2.3\ Repository\ shape\ at\ the\ pinned\ commit}
\item
  Purpose: Show the visible weight of the repo by top-level area and
  make the root map legible in one glance.
\item
  Source basis:
  \href{https://github.com/chris-page-gov/mcp-geo/blob/fe862910da246ca77f374cfbe484985f5df4d316/README.md}{\texttt{README.md}},
  \href{https://github.com/chris-page-gov/mcp-geo/blob/fe862910da246ca77f374cfbe484985f5df4d316/docs/reports/README.md}{\texttt{docs/reports/README.md}},
  \href{https://github.com/chris-page-gov/mcp-geo/blob/fe862910da246ca77f374cfbe484985f5df4d316/docs/public_sector_ai_community/README.md}{\texttt{docs/public\_sector\_ai\_community/README.md}},
  and the pinned commit root tree.
\item
  Intended output: \texttt{16:9} PNG suitable for PDF embedding.
\item
  Alt text draft: ``An infographic showing the top-level shape of the
  pinned MCP-Geo repository with the largest areas highlighted as docs,
  tests, research, scripts, server, tools, and UI.''
\item
  Prompt: Create a clean editorial infographic titled ``MCP-Geo
  repository at a glance''. Use an off-white paper background, dark ink
  text, muted teal and amber accents, and a public-sector research tone.
  Show the repository root as a central frame, then group the tracked
  top-level areas into meaningful clusters: root launch files, runtime
  code, tools and UI, documentation, research, testing and evidence,
  release operations, and support data. Include relative scale cues so
  docs, tests, and research read as the largest visible surfaces. Add
  small labels for representative anchors such as README.md, server/,
  tools/, docs/, tests/, research/, scripts/, and RELEASE\_NOTES/. Make
  the layout readable without surrounding text and suitable for a PDF
  report.
\end{itemize}

\subsection{F02. Reader pathways and start-here
matrix}\label{f02.-reader-pathways-and-start-here-matrix}

This prompt defines the decision aid that helps readers choose a route
through the repo based on purpose instead of guessing from folder names.

\begin{itemize}
\tightlist
\item
  Figure file: \texttt{f02\_reader\_pathways.png}
\item
  Placement: \texttt{\#\#\#\ 1.1\ Reading\ lanes\ and\ audience\ fit}
\item
  Purpose: Help readers choose the fastest route through the repo based
  on goal rather than folder name.
\item
  Source basis:
  \href{https://github.com/chris-page-gov/mcp-geo/blob/fe862910da246ca77f374cfbe484985f5df4d316/README.md}{\texttt{README.md}},
  \href{https://github.com/chris-page-gov/mcp-geo/blob/fe862910da246ca77f374cfbe484985f5df4d316/docs/public_sector_ai_community/README.md}{\texttt{docs/public\_sector\_ai\_community/README.md}},
  \href{https://github.com/chris-page-gov/mcp-geo/blob/fe862910da246ca77f374cfbe484985f5df4d316/docs/spec_package/README.md}{\texttt{docs/spec\_package/README.md}},
  and
  \href{https://github.com/chris-page-gov/mcp-geo/blob/fe862910da246ca77f374cfbe484985f5df4d316/docs/reports/README.md}{\texttt{docs/reports/README.md}}.
\item
  Intended output: \texttt{16:9} PNG suitable for PDF embedding.
\item
  Alt text draft: ``A matrix infographic mapping four reader goals to
  recommended repo reading paths, with arrows from onboarding, runtime,
  evidence, and publication needs to the best starting documents.''
\item
  Prompt: Create an infographic titled ``Reader pathways and start-here
  matrix'' for the MCP-Geo repository. Use four audience lanes:
  first-time operator, technical reviewer, evidence seeker, and
  publication editor. For each lane, show a short route through named
  repo documents such as README.md,
  docs/public\_sector\_ai\_community/README.md,
  docs/spec\_package/README.md, docs/reports/README.md, server/main.py,
  and server/mcp/tools.py. Make the layout work as a quick decision aid,
  with arrows and concise labels rather than prose-heavy paragraphs.
  Keep the style calm, professional, and clearly navigational.
\end{itemize}

\subsection{F03. Runtime request flow}\label{f03.-runtime-request-flow}

This prompt defines the transport-to-answer diagram that explains the
runtime stack to readers who need the system shape fast.

\begin{itemize}
\tightlist
\item
  Figure file: \texttt{f03\_runtime\_request\_flow.png}
\item
  Placement: \texttt{\#\#\ 3.\ Runtime\ and\ Protocol\ Implementation}
\item
  Purpose: Explain how a host request moves through transport, MCP
  routing, tools, resources, and observability.
\item
  Source basis:
  \href{https://github.com/chris-page-gov/mcp-geo/blob/fe862910da246ca77f374cfbe484985f5df4d316/server/main.py}{\texttt{server/main.py}},
  \href{https://github.com/chris-page-gov/mcp-geo/blob/fe862910da246ca77f374cfbe484985f5df4d316/server/stdio_adapter.py}{\texttt{server/stdio\_adapter.py}},
  \href{https://github.com/chris-page-gov/mcp-geo/blob/fe862910da246ca77f374cfbe484985f5df4d316/server/mcp/tools.py}{\texttt{server/mcp/tools.py}},
  \href{https://github.com/chris-page-gov/mcp-geo/blob/fe862910da246ca77f374cfbe484985f5df4d316/server/mcp/resources.py}{\texttt{server/mcp/resources.py}},
  and
  \href{https://github.com/chris-page-gov/mcp-geo/tree/fe862910da246ca77f374cfbe484985f5df4d316/server/audit}{\texttt{server/audit/}}.
\item
  Intended output: \texttt{16:9} PNG suitable for PDF embedding.
\item
  Alt text draft: ``A left-to-right runtime flow diagram showing host
  request entry, HTTP and STDIO transport, MCP routing, tool and
  resource dispatch, caches and audit, and response packaging.''
\item
  Prompt: Create a systems infographic titled ``Runtime request flow''.
  Show a left-to-right governed request path for MCP-Geo: host or
  client, HTTP or STDIO transport, MCP router, tool and resource
  dispatch, caches and helper subsystems, audit and observability, then
  the returned answer. Include representative file anchors such as
  server/main.py, server/stdio\_adapter.py, server/mcp/tools.py,
  server/mcp/resources.py, and server/audit/. Make the visual
  understandable to non-specialist readers while still feeling
  technically credible. Use clear arrows, a small legend, and restrained
  colors.
\end{itemize}

\subsection{F04. Tool and resource ecosystem
map}\label{f04.-tool-and-resource-ecosystem-map}

This prompt defines the capability map that lets readers understand the
callable surface before opening the schema catalog.

\begin{itemize}
\tightlist
\item
  Figure file: \texttt{f04\_tool\_resource\_ecosystem.png}
\item
  Placement:
  \texttt{\#\#\#\ 4.2\ Domain\ tool\ families\ and\ supporting\ resources}
\item
  Purpose: Show the breadth of callable tool families without forcing
  readers through the full schema catalog first.
\item
  Source basis:
  \href{https://github.com/chris-page-gov/mcp-geo/blob/fe862910da246ca77f374cfbe484985f5df4d316/docs/tool_catalog.md\#L1-L61}{\texttt{docs/tool\_catalog.md}},
  \href{https://github.com/chris-page-gov/mcp-geo/blob/fe862910da246ca77f374cfbe484985f5df4d316/server/mcp/tools.py}{\texttt{server/mcp/tools.py}},
  \href{https://github.com/chris-page-gov/mcp-geo/tree/fe862910da246ca77f374cfbe484985f5df4d316/resources}{\texttt{resources/}},
  and
  \href{https://github.com/chris-page-gov/mcp-geo/tree/fe862910da246ca77f374cfbe484985f5df4d316/ui}{\texttt{ui/}}.
\item
  Intended output: \texttt{16:9} PNG suitable for PDF embedding.
\item
  Alt text draft: ``A hub-and-spoke diagram showing MCP-Geo tool
  families around a central registry, with related resources and UI
  surfaces connected as supporting layers.''
\item
  Prompt: Create an infographic titled ``Tool and resource ecosystem
  map'' for MCP-Geo. Put the tool registry at the center, then arrange
  the main tool families around it: OS Places and Names, OS features and
  maps, admin lookup, ONS and NOMIS, route planning, MCP meta tools, and
  MCP-Apps UI tools. Add a secondary ring for resources and UI assets,
  including resources/, docs/tool\_catalog.md, ui/, and playground/. Use
  concise labels and category colors so a visual learner can grasp the
  breadth of the server without reading the full schema catalog. Keep
  the visual polished and publication-ready.
\end{itemize}

\subsection{F05. Documentation and publication
stack}\label{f05.-documentation-and-publication-stack}

This prompt defines the document-stack visual that keeps the explanatory
surfaces legible instead of burying them inside \texttt{docs/}.

\begin{itemize}
\tightlist
\item
  Figure file: \texttt{f05\_documentation\_publication\_stack.png}
\item
  Placement:
  \texttt{\#\#\ 5.\ Documentation,\ Reports,\ Slides,\ and\ Publication\ Paths}
\item
  Purpose: Show how onboarding notes, specs, reports, slides, and Prism
  outputs relate instead of competing with each other.
\item
  Source basis:
  \href{https://github.com/chris-page-gov/mcp-geo/blob/fe862910da246ca77f374cfbe484985f5df4d316/docs/getting_started.md}{\texttt{docs/getting\_started.md}},
  \href{https://github.com/chris-page-gov/mcp-geo/blob/fe862910da246ca77f374cfbe484985f5df4d316/docs/spec_package/README.md}{\texttt{docs/spec\_package/README.md}},
  \href{https://github.com/chris-page-gov/mcp-geo/blob/fe862910da246ca77f374cfbe484985f5df4d316/docs/public_sector_ai_community/README.md}{\texttt{docs/public\_sector\_ai\_community/README.md}},
  \href{https://github.com/chris-page-gov/mcp-geo/blob/fe862910da246ca77f374cfbe484985f5df4d316/docs/reports/README.md}{\texttt{docs/reports/README.md}},
  and
  \href{https://github.com/chris-page-gov/mcp-geo/tree/fe862910da246ca77f374cfbe484985f5df4d316/docs/mcp_geo_prism_bundle}{\texttt{docs/mcp\_geo\_prism\_bundle/}}.
\item
  Intended output: \texttt{16:9} PNG suitable for PDF embedding.
\item
  Alt text draft: ``A stacked publication diagram showing onboarding
  docs, specifications, narrative documentation, reports, slides, and
  Prism bundle outputs as linked layers.''
\item
  Prompt: Create an infographic titled ``Documentation and publication
  stack''. Show MCP-Geo's explanatory material as a layered stack:
  quick-start docs, specification package, community narrative set,
  report library, slide and event materials, then Prism-ready
  publication output. Use representative labels such as
  docs/getting\_started.md, docs/spec\_package/README.md,
  docs/public\_sector\_ai\_community/README.md, docs/reports/README.md,
  docs/slides/20260225 - From\_Apps\_to\_Answers.pdf, and
  docs/mcp\_geo\_prism\_bundle/. Make the layout feel like a publishing
  pipeline rather than a generic folder chart.
\end{itemize}

\subsection{F06. Evidence and assurance
ladder}\label{f06.-evidence-and-assurance-ladder}

This prompt defines the ladder diagram that shows how MCP-Geo supports
scrutiny instead of only feature discovery.

\begin{itemize}
\tightlist
\item
  Figure file: \texttt{f06\_evidence\_assurance\_ladder.png}
\item
  Placement:
  \texttt{\#\#\#\ 6.2\ Tests,\ evaluation\ harnesses,\ and\ benchmark\ packs}
\item
  Purpose: Show how research, tests, evaluation, reports, and audit
  materials combine into a defensible evidence trail.
\item
  Source basis:
  \href{https://github.com/chris-page-gov/mcp-geo/tree/fe862910da246ca77f374cfbe484985f5df4d316/tests}{\texttt{tests/}},
  \href{https://github.com/chris-page-gov/mcp-geo/tree/fe862910da246ca77f374cfbe484985f5df4d316/tests/evaluation}{\texttt{tests/evaluation/}},
  \href{https://github.com/chris-page-gov/mcp-geo/tree/fe862910da246ca77f374cfbe484985f5df4d316/docs/reports}{\texttt{docs/reports/}},
  \href{https://github.com/chris-page-gov/mcp-geo/blob/fe862910da246ca77f374cfbe484985f5df4d316/docs/decision_support_audit_pack.md}{\texttt{docs/decision\_support\_audit\_pack.md}},
  and
  \href{https://github.com/chris-page-gov/mcp-geo/tree/fe862910da246ca77f374cfbe484985f5df4d316/troubleshooting}{\texttt{troubleshooting/}}.
\item
  Intended output: \texttt{16:9} PNG suitable for PDF embedding.
\item
  Alt text draft: ``A ladder diagram showing successive assurance layers
  from raw research and fixtures through tests, evaluation harnesses,
  reports, troubleshooting, and audit packaging.''
\item
  Prompt: Create an infographic titled ``Evidence and assurance
  ladder''. Present MCP-Geo's assurance model as ascending steps:
  research inputs and option papers, fixtures and scripts, automated
  tests, evaluation harnesses and benchmark packs, human-readable
  reports, troubleshooting traces, and audit-ready packaging. Use
  repository anchors such as research/, tests/, tests/evaluation/,
  docs/reports/, troubleshooting/, and
  docs/decision\_support\_audit\_pack.md. The visual should explain
  rigor and traceability to a reader who may never open the code.
\end{itemize}

\subsection{F07. Release and iteration
pipeline}\label{f07.-release-and-iteration-pipeline}

This prompt defines the publication-and-iteration loop that turns repo
maintenance into a visible maturity signal.

\begin{itemize}
\tightlist
\item
  Figure file: \texttt{f07\_release\_iteration\_pipeline.png}
\item
  Placement:
  \texttt{\#\#\#\ 6.3\ Release,\ CI,\ troubleshooting,\ and\ maturity\ signals}
\item
  Purpose: Show how work moves from implementation to verification to
  public release and then back into documentation.
\item
  Source basis:
  \href{https://github.com/chris-page-gov/mcp-geo/blob/fe862910da246ca77f374cfbe484985f5df4d316/CHANGELOG.md}{\texttt{CHANGELOG.md}},
  \href{https://github.com/chris-page-gov/mcp-geo/tree/fe862910da246ca77f374cfbe484985f5df4d316/RELEASE_NOTES}{\texttt{RELEASE\_NOTES/}},
  \href{https://github.com/chris-page-gov/mcp-geo/blob/fe862910da246ca77f374cfbe484985f5df4d316/.github/workflows/ci.yml}{\texttt{.github/workflows/ci.yml}},
  \href{https://github.com/chris-page-gov/mcp-geo/blob/fe862910da246ca77f374cfbe484985f5df4d316/PROGRESS.MD}{\texttt{PROGRESS.MD}},
  and
  \href{https://github.com/chris-page-gov/mcp-geo/blob/fe862910da246ca77f374cfbe484985f5df4d316/CONTEXT.md}{\texttt{CONTEXT.md}}.
\item
  Intended output: \texttt{16:9} PNG suitable for PDF embedding.
\item
  Alt text draft: ``A pipeline diagram showing feature delivery, tests
  and CI, release packaging, documentation refresh, and feedback loops
  back into context and progress tracking.''
\item
  Prompt: Create an infographic titled ``Release and iteration
  pipeline''. Show a clear loop: implementation work, scripts and
  tooling, automated tests and CI, release notes and changelog, public
  publication, then context and progress updates feeding the next
  iteration. Label repository anchors such as .github/workflows/ci.yml,
  CHANGELOG.md, RELEASE\_NOTES/, PROGRESS.MD, and CONTEXT.md. Keep the
  flow explicit and governance-oriented, suitable for a UK public-sector
  AI discussion.
\end{itemize}

\subsection{F08. Governed answer loop}\label{f08.-governed-answer-loop}

This prompt defines the concept figure that reconnects the repo
structure to the original talk's core argument.

\begin{itemize}
\tightlist
\item
  Figure file: \texttt{f08\_governed\_answer\_loop.png}
\item
  Placement: \texttt{\#\#\#\ 5.3\ Prism\ and\ publication\ stack}
\item
  Purpose: Tie the talk thesis to the repo by showing how a question
  becomes an inspectable answer and then a publication artifact.
\item
  Source basis:
  \href{https://github.com/chris-page-gov/mcp-geo/blob/fe862910da246ca77f374cfbe484985f5df4d316/docs/slides/20260225\%20-\%20From_Apps_to_Answers.pdf}{\texttt{docs/slides/20260225\ -\ From\_Apps\_to\_Answers.pdf}},
  \href{https://github.com/chris-page-gov/mcp-geo/blob/fe862910da246ca77f374cfbe484985f5df4d316/docs/reports/mcp_geo_ai_community_event_record.docx}{\texttt{docs/reports/mcp\_geo\_ai\_community\_event\_record.docx}},
  \href{https://github.com/chris-page-gov/mcp-geo/tree/fe862910da246ca77f374cfbe484985f5df4d316/server/audit}{\texttt{server/audit/}},
  and
  \href{https://github.com/chris-page-gov/mcp-geo/tree/fe862910da246ca77f374cfbe484985f5df4d316/docs/mcp_geo_prism_bundle}{\texttt{docs/mcp\_geo\_prism\_bundle/}}.
\item
  Intended output: \texttt{16:9} PNG suitable for PDF embedding.
\item
  Alt text draft: ``A circular flow diagram showing a question entering
  MCP-Geo, passing through tools and evidence, producing an answer and
  audit trail, and being turned into report and publication outputs.''
\item
  Prompt: Create an infographic titled ``Governed answer loop''. Show a
  circular public-sector evidence workflow: question, MCP routing and
  tool choice, evidence retrieval, answer assembly, provenance and audit
  capture, human-readable report, Prism or PDF publication, then
  re-entry into future questions. Include light references to the talk
  material and event record without making the slide deck the whole
  story. The visual should make the repo feel like a governed answer
  system rather than a loose collection of files.
\end{itemize}
